\newpage\thispagestyle{hl}
\hlbanner{Part III $\cdot$ Reduced-order modelling}{11}{Substructuring and component reduction}
\hlsection{The question}
\begin{hlquote}A wing, an engine and a pylon come from different teams, and each team can reduce its own part. The difficulty is joining the reduced parts again: if the reduction hides the interface the parts cannot be connected, and if it keeps every interface unknown the interface dominates.\end{hlquote}
\vspace{8pt}\begin{hlsum}\textbf{Each part is reduced while its interface unknowns are kept, as Craig--Bampton does.} For a clamped--clamped beam of \ccmBeamFull{} unknowns split in two, Craig--Bampton with \ccmCbFiveSize{} unknowns has an error of \ccmCbFiveErr, while Guyan condensation at the same size errs by \ccmGuyTwelveErr, about twelve times more. The chapter explains both, judges an assembly by its frequencies, and says what free-interface and interface reduction change. The tests are small (a beam, a membrane, a spring chain), with no solid, no shell and no assembly of more than two parts.\end{hlsum}
\hlsection{Key concepts}
\begin{description}\setlength\itemsep{5pt}\setlength\parsep{0pt}\renewcommand{\descriptionlabel}[1]{\hspace{\labelsep}\sffamily\bfseries\color{hlnavy}#1}
\item[Static (Guyan) condensation] eliminates interior unknowns statically; exact for statics, weak in dynamics
\item[Craig--Bampton] constraint modes plus fixed-interface modes, with a physical interface
\item[Free-interface reduction] free modes with residual flexibility of the discarded ones
\item[Interface reduction] a few interface shapes replace every interface unknown
\item[Component POD, mortar tie] component-wise snapshots joined across non-matching meshes
\item[Judging an assembly] by its frequencies against the unreduced model
\end{description}
\hlsection{After this chapter you can}
\begin{itemize}\setlength\itemsep{3pt}
\item build and compare Guyan and Craig--Bampton models
\item judge an assembly by its frequencies
\item say what free-interface and interface reduction change
\end{itemize}
\newpage\thispagestyle{hl}
\noindent{\sffamily\small\color{hlgrey}Chapter 11 \quad Substructuring and component reduction}\par
\hlfig{\hlimg{ch11_cb}}{Largest error of the first four frequencies of a two-part beam against reduced model size, for Craig--Bampton and Guyan condensation.}
\hlsection{Headline result}
\begin{hlheadline}For a clamped--clamped beam of \ccmBeamFull{} unknowns split in two, Craig--Bampton with \ccmCbFiveSize{} unknowns has an error of \ccmCbFiveErr. Guyan condensation at the same size errs by \ccmGuyTwelveErr, about a factor 12 larger.\end{hlheadline}
\hlsection{Main classes and functions}
{\begin{itemize}\setlength\itemsep{2pt}
\item \texttt{rom\_engine.cms.guyan}
\item \texttt{rom\_engine.interface\_reduction.interface\_reduced\_craig\_bampton}
\item \texttt{rom\_engine.component\_pod.ReducedAssembly}
\item \texttt{rom\_engine.component\_pod.ComponentBasis}
\end{itemize}}
\hlsection{A stated limit}
\begin{hllimit}Small, simple test models: a beam, a membrane and a spring chain. There is no three-dimensional solid substructure, no shell and no assembly of more than two parts.\end{hllimit}
